\section{Proposed method}\label{sec:method}
\mname{} (lead-time-scaled chance constraints) takes four decisions in every replanning cycle
(Fig.~\ref{fig:overview}, Algorithm~\ref{alg:cycle}). The chance-constrained planning model decides which plans are
acceptable under the calibrated execution error. The cost-aware fallback decides which of the tasks that this model
gives up are still worth attempting, given the cost of a late arrival. The stopping rule decides when to stop computing
and publish, and the trigger decides when to replan. The first two decisions are the contribution of this paper. Both
reduce to one effective deadline per task, so a feasibility-first repair and an anytime ALNS, or any other heuristic
that checks deadlines, produce the plan unchanged. Table~\ref{tab:notation} lists the notation. Two earlier variants
serve as baselines. AnytimeSwarm uses the same stopping rule with a deterministic delay cost on the nominal model, and
AnytimeSwarm-R uses an execution-aware delay cost on the nominal model.

\begin{table}[t]
\caption{Notation.}\label{tab:notation}
\centering\footnotesize
\setlength{\tabcolsep}{3pt}
\begin{tabular}{@{}ll@{}}
\toprule
Symbol & Meaning \\
\midrule
$t_0$ & time of the planning snapshot \\
$r_j,\,d_j,\,w_j$ & release time, deadline and reward of task $j$ \\
$f_j,\,F_j$ & planned (nominal) and actual finish time of task $j$ \\
$h_j=f_j-t_0$ & lead time of task $j$ \\
$e_j=F_j-f_j$ & finish-time error, $e_j\sim\mathcal{N}(m(h_j),s(h_j)^2)$ \\
$m(h),\,s(h)$ & mean and spread of the error, $b_0+b_1h$ and $s_0+s_1h$ \\
$z$ & safety factor; protected tasks meet $d_j$ with probability $\ge\Phi(z)$ \\
$\tilde d_j(z)$ & effective deadline, Eq.~\eqref{eq:cc} \\
$\tilde H,\,H$ & tasks no UAV can finish by $\tilde d_j$, resp.\ by $d_j$, on a direct flight \\
$\lambda$ & cost of a late arrival, in units of the task reward \\
$\theta(\lambda),\,z_\lambda$ & break-even probability $\lambda/(1+\lambda)$ and $\Phi^{-1}(\theta(\lambda))$ \\
$d'_j$ & planning deadline used by repair and search \\
$\tau,\,B$ & planning time of a call and its hard budget \\
$\kappa,\,L$ & latency factor and communication latency \\
$u_\lambda$ & score $U-\lambda\,n_{\mathrm{late}}/n$ of a mission \\
\bottomrule
\end{tabular}
\end{table}

\begin{figure*}[t]
\centering
\fig{fig_overview}{\linewidth}
\caption{Overview of \mname{}. Top: the replanning loop between the mission, the planning snapshot, the
replanner and the execution tier. Bottom: one replanning call. The calibrated error model gives a margin that grows
with the lead time, the chance constraint turns it into an effective deadline per task, the cost-aware fallback keeps a
given-up task only if its probability of meeting the deadline outweighs the cost $\lambda$ of a late arrival, and the
anytime search stops when the predicted gain no longer pays for the delay. The plan is released after the emulated planning delay $\kappa\tau$ and the latency $L$.}\label{fig:overview}
\end{figure*}

\begin{algorithm}[t]
\caption{One replanning call of \mname{}}\label{alg:cycle}
\small
\begin{algorithmic}[1]
\Require snapshot at time $t_0$: committed tasks, open tasks with $(r_j,d_j,w_j)$, UAV states; previous plan $\pi^-$;
error model $(b_0,b_1,s_0,s_1)$, safety factor $z$, late cost $\lambda$, expected latency $\bar L$, latency factor
$\kappa$
\State $z_\lambda \gets \Phi^{-1}(\lambda/(1+\lambda))$;\; $z_p\gets\max\{z,z_\lambda\}$
\Comment{$z_\lambda=-\infty$ for $\lambda=0$}
\State $\tilde d_j(\cdot) \gets \min\{d_j,\; t_0+(d_j-t_0-b_0-\cdot\, s_0)/(1+b_1+\cdot\, s_1)\}$ for every open task
\Comment{Eq.~\eqref{eq:cc}}
\State $d'_j \gets \min\{d_j,\tilde d_j(z_\lambda)\}$ if no UAV can finish $j$ by $\tilde d_j(z_p)$ but one can by
$d_j$, else $d'_j\gets\tilde d_j(z_p)$ \Comment{Eq.~\eqref{eq:costfb}}
\State $\pi \gets$ \Call{FeasibilityFirstRepair}{$\pi^-$, $d'$};\; $\pi^\ast \gets \pi$;\; $\tau \gets 0$
\Repeat
  \State destroy and repair $\pi$ with adaptively chosen operators; accept by simulated annealing
  \State \textbf{if} $J(\pi) > J(\pi^\ast)$ \textbf{then} $\pi^\ast \gets \pi$
  \State every 5\,ms: $\widehat{\Delta J} \gets$ learned prediction of the gain in the next $H$\,ms
\Until{$\widehat{\Delta J}\le C_{\mathrm{R}}(\tau)+\mu W H$ \textbf{or} $\tau\ge B$} \Comment{Eq.~\eqref{eq:stop}}
\State publish $\pi^\ast$; it reaches the UAVs at $t_0+\kappa\tau+L$
\State next call at the next event or $T_{\max}$ after $t_0$, whichever comes first
\end{algorithmic}
\end{algorithm}

\subsection{Repair and anytime search}
The previous plan is restricted to the open tasks and walked once, keeping a task only if the route up to it is still
feasible, and tasks of failed UAVs and new tasks are inserted by parallel regret-2
insertion~\cite{ropke2006adaptive}. This feasibility-first repair yields a publishable plan within
milliseconds. ALNS~\cite{ropke2006adaptive} then improves it with four destroy operators (random, worst,
related and route-segment removal), four greedy or regret-2 repair operators with and without noise,
adaptive operator selection and simulated-annealing acceptance, the standard components of ALNS~\cite{ropke2006adaptive,mara2022alns}. All intermediate
solutions are feasible, so the best plan found so far can be published at any time. The temperature decreases per
iteration and all variants share the random seed, so every stopping rule follows the same search trajectory and differs
only in the moment at which it stops.

\subsection{Calibrated finish-time error}
The error of task $j$ is modelled as $e_j\sim\mathcal{N}(m(h_j),s(h_j)^2)$, where $h_j=f_j-t_0$ is its lead time, the
nominal finish time minus the time $t_0$ of the planning snapshot, and
\begin{equation}\label{eq:errmodel}
m(h)=b_0+b_1h,\qquad s(h)=s_0+s_1h ,
\end{equation}
because errors accumulate along a route. The parameters are fitted per disturbance level on development missions, the
mean by least squares and the spread by least squares on the absolute residuals scaled by $\sqrt{\pi/2}$, which is less
sensitive to the heavy tail of the errors. Under the mid and high disturbances of the main studies, both the mean and
the spread grow with the lead time (Supplementary Table~\ref{S-tab:errmodels}). The nominal model is therefore optimistic, and
increasingly so along a route.

The affine spread is a modelling choice that the data can reject. A route-level model with independent increments per
unit of time, the usual assumption when variances are added along a route, implies $s(h)=\sqrt{q_0+q_1h}$. On
\val{ENPairsMid} and \val{ENPairsHigh} pairs of lead time and error from \val{ENMissions} development missions per
level, the affine spread fits better by \val{EDAicSqrtMid} and \val{EDAicSqrtHigh} points of the Akaike information
criterion under mid and high disturbances, and also better than the hybrid $\sqrt{q_0+q_1h+q_2h^2}$ (by
\val{EDAicHybMid} and \val{EDAicHybHigh}). A five-fold cross-validation with missions as groups gives the same ranking.
At the longest lead times, about \val{EHLateHigh}\,s, the square-root model predicts a spread of \val{ESdLateSqrtHigh}\,s
under high disturbances where \val{ESdLateHigh}\,s are observed, while the affine model predicts \val{ESdLateAffHigh}\,s
(Fig.~\ref{fig:errstructure}a).

The reason is that the disturbances persist over a leg. Wind and an airspeed error act during the whole flight to a
task, so the error of a leg is proportional to its duration (Fig.~\ref{fig:errstructure}b), with a coefficient of
variation of \val{EMpCvMid} and \val{EMpCvHigh} under mid and high disturbances. Errors of successive legs are close to
independent. The correlation between the error of a task and the error added by the next task of the same plan was
\val{EIncCorrMid} and \val{EIncCorrHigh}, and missions explained only \val{EShareMissionMid}\% and
\val{EShareMissionHigh}\% of the residual variance. A lead time of a few legs therefore carries a spread that grows
almost linearly. Moment propagation with a variance proportional to the squared duration of each leg is consistent
with this structure and serves as the literature baseline in Section~\ref{sec:methods}.

\begin{figure*}[t]
\centering
\fig{fig_errstructure}{0.8\linewidth}
\caption{Structure of the finish-time error on development missions. (a) Standard deviation of the error against the
lead time, in deciles, with the affine fit of Eq.~\eqref{eq:errmodel} and the fit with independent increments. (b)
Standard deviation of the flight-time error of single legs against their nominal flight time, in deciles, with a
spread proportional to the duration.}\label{fig:errstructure}
\end{figure*}

\subsection{Chance-constrained planning model}\label{sec:cc}
The proposed model plans for the calibrated execution. A planned finish time $f_j$ is accepted only if
\begin{equation}\label{eq:ccfinish}
f_j+m(f_j-t_0)+z\,s(f_j-t_0)\;\le\;d_j ,
\end{equation}
where $z\ge0$ is a safety factor. The margin $m+zs$ depends on the finish time itself, so it is small for tasks served
soon and large for tasks served late in a route. Because $m$ and $s$ are linear, the constraint reduces to an effective
deadline,
\begin{equation}\label{eq:cc}
f_j\le\tilde d_j,\qquad \tilde d_j=t_0+\frac{d_j-t_0-b_0-z s_0}{1+b_1+z s_1}.
\end{equation}
The repair and the search use $\tilde d_j$ instead of $d_j$, while the executor and the reward still use the true
deadline.

\begin{prop}\label{prop:cc}
Let $1+b_1+z s_1>0$. (i) A finish time satisfies Eq.~\eqref{eq:ccfinish} if and only if $f_j\le\tilde d_j$. (ii) Under
the error model of Eq.~\eqref{eq:errmodel}, a task planned to finish at $f_j\le\tilde d_j$ meets its deadline with
probability at least $\Phi(z)$ if its route starts as planned. (iii) With $b_0=b_1=s_0=s_1=0$, $\tilde d_j=d_j$ and the
planning model is the nominal one.
\end{prop}
\begin{proof}
Substituting Eq.~\eqref{eq:errmodel} into Eq.~\eqref{eq:ccfinish} gives
$f_j(1+b_1+zs_1)\le d_j-b_0-zs_0+(b_1+zs_1)t_0$, and dividing by the positive factor gives (i). For (ii), the actual
finish is $F_j=f_j+e_j$ with $e_j\sim\mathcal{N}(m(h_j),s(h_j)^2)$ and $h_j=f_j-t_0$. By (i),
$d_j-f_j-m(h_j)\ge z\,s(h_j)$, hence $\Pr[F_j\le d_j]=\Phi\big((d_j-f_j-m(h_j))/s(h_j)\big)\ge\Phi(z)$. Part (iii) is
immediate.
\end{proof}

The effective deadline turns a chance constraint on every task into a deterministic constraint of the same form as the
original one. The insertion heuristics and the feasibility checks of the search are therefore reused unchanged, and
computing $\tilde d_j$ costs one division per task. With $z=0$, the model corrects the mean error only and plans for the
median finish time. Larger values of $z$ protect more tasks against lateness but leave fewer tasks feasible, so $z$
trades the risk of a late service against the reward of attempting a task. It is chosen per disturbance level on
development missions.

\subsection{Nominal fallback for given-up tasks}\label{sec:fallback}
A hard chance constraint has a side effect when the time windows are short compared with the spread of the error. If no
vehicle can finish task $j$ before $\tilde d_j$, even on a direct flight, the task is infeasible for the planner and is
never attempted. A vehicle may still serve it in time with a probability below $\Phi(z)$ but well above zero, and it may
have nothing better to do. The fallback keeps such tasks. Let $\tilde H$ be the set of tasks that no vehicle can finish
by $\tilde d_j$ on a direct flight, and $H$ the set of tasks that no vehicle can finish by $d_j$. The planning deadline is
\begin{equation}\label{eq:fallback}
d'_j=\begin{cases} d_j, & j\in\tilde H\setminus H,\\ \tilde d_j, & \text{otherwise.}\end{cases}
\end{equation}
The constraints of all other tasks are unchanged, so every other task that is served keeps the guarantee of
Proposition~\ref{prop:cc}, while a fallback task carries none. Finding $\tilde H$ and $H$ requires one direct-flight
check per task and vehicle, which is negligible compared with the search. The fallback was designed after the held-out
results of study~2 and is evaluated on new seeds in studies~3 and~4.

\subsection{Cost-aware fallback}\label{sec:costfb}
The nominal fallback is right only if a late service costs nothing. In many applications a late arrival has a cost of
its own. The flight time and energy are spent, a customer may have left, or a commitment is broken. Let a late arrival
at task $j$ cost $\lambda w_j$, so that $\lambda$ measures this cost in units of the reward of the task. An attempt that
meets the deadline with probability $p$ then has the expected value $p\,w_j-(1-p)\lambda w_j$, which is non-negative if
and only if
\begin{equation}\label{eq:theta}
p\;\ge\;\theta(\lambda)=\frac{\lambda}{1+\lambda}.
\end{equation}
By Proposition~\ref{prop:cc}, a finish time $f_j\le\tilde d_j(z)$ gives a probability of at least $\Phi(z)$, where
$\tilde d_j(z)$ denotes the effective deadline of Eq.~\eqref{eq:cc} at safety factor $z$. The break-even probability
therefore corresponds to the safety factor $z_\lambda=\Phi^{-1}(\theta(\lambda))$. The cost-aware fallback protects
every task at $z_p=\max(z,z_\lambda)$, and plans a task that no vehicle can finish by $\tilde d_j(z_p)$ on a direct
flight against $\min\{d_j,\tilde d_j(z_\lambda)\}$. A task that cannot meet this deadline either is given up. Thus
\begin{equation}\label{eq:costfb}
d'_j=\begin{cases} \min\{d_j,\tilde d_j(z_\lambda)\}, & j\in\tilde H(z_p)\setminus H,\\ \tilde d_j(z_p), &
\text{otherwise,}\end{cases}
\end{equation}
with $\tilde H(z_p)$ the set of tasks that no vehicle can finish by $\tilde d_j(z_p)$ on a direct flight.

\begin{prop}\label{prop:costfb}
Let $s_0,s_1\ge0$, let $1+b_1+z s_1>0$ for every finite $z$ used, and read $\tilde d_j(-\infty)$ as $d_j$. (i) For
$\lambda=0$, Eq.~\eqref{eq:costfb} is the nominal fallback of Eq.~\eqref{eq:fallback}. (ii) For
$\lambda\ge\Phi(z)/(1-\Phi(z))$, no task falls back, and every planned task meets $\tilde d_j(z_\lambda)$. (iii) Under
the error model of Eq.~\eqref{eq:errmodel}, every task planned to finish by $d'_j$ has an expected contribution
$w_j\Pr[F_j\le d_j]-\lambda w_j\Pr[F_j>d_j]\ge0$ if its route starts as planned. (iv) Every planning deadline $d'_j$ is
non-increasing in $\lambda$.
\end{prop}
\begin{proof}
(i) $z_\lambda=-\infty$ gives $z_p=z$ and $\min\{d_j,\tilde d_j(-\infty)\}=d_j$. (ii) $z_\lambda\ge z$ holds if and only
if $\theta(\lambda)\ge\Phi(z)$, that is, $\lambda\ge\Phi(z)/(1-\Phi(z))$. Then $z_p=z_\lambda$, and a task that cannot
meet $\tilde d_j(z_p)$ cannot meet $\min\{d_j,\tilde d_j(z_\lambda)\}$ either, so it is given up. (iii) A protected task
meets its deadline with probability at least $\Phi(z_p)\ge\Phi(z_\lambda)$ and a fallback task with probability at least
$\Phi(z_\lambda)$, both by Proposition~\ref{prop:cc}. Hence $p\ge\theta(\lambda)$, and the expected contribution
$w_j\big((1+\lambda)p-\lambda\big)$ is not negative. (iv) With $s_0,s_1\ge0$, $\tilde d_j(z)$ does not increase with
$z$, and $z_\lambda$ and $z_p$ do not decrease with $\lambda$. A fallback task stays one when $\lambda$ grows, because
$z_p$ does not decrease. A protected task can become a fallback task only while $z_\lambda<z$, and then $z_p=z$ does not
change. In each case the deadline does not increase.
\end{proof}

For $\lambda$ between these limits, the rule attempts a given-up task only if its expected contribution is not
negative. It adds no parameter to tune, because $z$ keeps its value and $\lambda$ is a property of the application.
Where $1+b_1+z_\lambda s_1\le0.05$, which occurs only for $\lambda<0.03$, the nominal deadline is used. The rule needs
only a probability of meeting the deadline, so it applies to any planning model that provides one, such as moment
propagation, where the given-up tasks are accepted if $\mathrm{E}[F_j]+z_\lambda\mathrm{sd}[F_j]\le d_j$. It ignores
the delay that a fallback task causes to the tasks after it on the same route. A margin without a probability of
meeting the deadline, such as a proportional buffer, has no counterpart of $z_\lambda$, so its fallback deadline has to
be tuned for every value of $\lambda$. The cost-aware fallback was designed after the held-out results of study~4 and is
evaluated on new seeds in studies~5 and~6.

\subsection{Stopping rule and trigger}\label{sec:stop}
Every 5\,ms, the stopping rule compares the predicted gain $\widehat{\Delta J}$ of searching $H=10$\,ms longer with the
cost of doing so, and stops at time $\tau$ if
\begin{equation}\label{eq:stop}
\widehat{\Delta J}\;\le\;C_{\mathrm{R}}(\tau)+\mu W H ,
\end{equation}
with a price of computation $\mu=10^{-5}$ per millisecond. The gain is predicted from eleven features of the search
history by a Poisson generalised linear model~\cite{mccullagh1989glm} fitted on development trajectories. The delay
cost $C_{\mathrm{R}}(\tau)=\Lambda_{\mathrm{R}}(\tau+H)-\Lambda_{\mathrm{R}}(\tau)$ is the increase of the expected value
of the tasks that a later start would lose,
\begin{equation}\label{eq:risk}
\Lambda_{\mathrm{R}}(\tau)=\sum_j v_j\,\Phi\!\left(\frac{\sigma_j(\tau)+\bar L+m(h_j)-\mathrm{tol}_j}{s(h_j)}\right),
\end{equation}
where $\sigma_j(\tau)$ is the start shift of the route of task $j$ caused by waiting for the plan, $\mathrm{tol}_j$ the
margin of the task to its deadline plus the waiting times before it on its route~\cite{savelsbergh1992vrptw}, $\bar L$
the expected latency and $v_j$ the value of the task in Eq.~\eqref{eq:obj}. On the effective deadlines, the mean error
is already part of the tolerance, so \mname{} uses $m\equiv0$. $\Lambda_{\mathrm{R}}$ does not decrease with
$\tau$, and it reduces to the exact loss of a deterministic plan when the spread vanishes (Supplementary
Section~\ref{S-sec:Sstop}). A hard budget $B$ bounds every call. The replanner replans on every event and at least every
$T_{\max}=60$\,s of mission time, a rolling-horizon policy in the terminology of
dynamic routing~\cite{soeffker2022review}.
